\documentclass[a4paper,10pt]{article}
% ali \documentclass{beamer}
\usepackage{babel}
\usepackage[utf8]{inputenc}
\usepackage[T1]{fontenc} %prvi trije nujni za slovenščino
\usepackage{amsmath} %poravnane enačbe, matrike ipd.
\usepackage{amsthm} %definicije, trditve
\usepackage{amssymb} %matematični simboli
\usepackage{graphicx} %slike(\inlcudegraphics)
\usepackage{booktabs} %lepe tabele
\usepackage{tikz} %slike/grafi
\usepackage{pgfplots} %grafi
\usepackage{blindtext} %za blind text (ni pomembno)
\usepackage{hyperref} %povezave (naj bo zadnji)

\newcommand{\R}{\mathbb{R}}



% Definicija okolij izrek, posledica
{\theoremstyle{plain}
\newtheorem{izrek}{Izrek}[section]
\newtheorem{posledica}[izrek]{Posledica}
}

% Definicija okoli za definicije in vaje
{\theoremstyle{definition}
\newtheorem{definicija}[izrek]{Definicija}
\newtheorem{vaja}[izrek]{Vaja}
}

%Lahko tudi samo tako:
%\theoremstyle{definition}
%\newtheorem{definicija}{Definicija}

%\theoremstyle{plain}
%\newtheorem{trditev}{Trditev}

\title{Naslov}
\author{Ime}
\date{\today} %oziroma ce noces datuma: \date{}
\begin{document}
\maketitle

\begin{abstract}
Uporabljamo za povzetek.
\end{abstract}

\tableofcontents %za seznam poglavij

\section{
    To je prvo poglavje
}
\blindtext
\subsection{
    To je podpoglavje prvega poglavja
}
\blindtext

%\input{...} za dodajanje datotek

\emph{poudarjevo}, \textbf{krepko}, \textit{poševno}, \texttt{pisava za kodo},
 \verb|\maketitle|, %za prikaz kode
  \url{https://youtube.com}, \noindent, 
  \centering Blabla
 \begin{verbatim}            dobesedno besedilo      sdsdsdsd
 \end{verbatim}

\begin{itemize}
  \item Prvi
  \item Drugi
\end{itemize}

\begin{enumerate}
  \item Ena
  \item Dva
\end{enumerate}

\begin{table}[h]
\centering
\caption{Primer tabele}
\label{tab:rezultati}
\begin{tabular}{lcr}
\toprule
\textbf{Ime} & \textbf{Točke} & \textbf{Ocena} \\
\midrule
Ana   & 92 & 10 \\
Boris & 78 & 8  \\
Cene  & 64 & 7  \\
\midrule
\multicolumn{3}{c}{Povprečje: 78} \\
\bottomrule
\end{tabular}
\end{table} %ta nima navpicnih crt


\begin{table}[ht]
\centering
\caption{Rezultati izpita pri predmetu Računalniški praktikum}
\label{tab:izpit}
\begin{tabular}{|l|c|c|c|r|} %|za navpicne crte
\hline
\textbf{Ime} & \textbf{Vpisna št.} & \textbf{Naloga 1} & \textbf{Naloga 2} & \textbf{Skupaj} \\
\hline
Ana Novak   & 631201 & 18 & 20 & 38 \\ %\\za prelom vrstice
Boris Kranj & 631202 & 15 & 17 & 32 \\
Cene Horvat & 631203 & 12 & 14 & 26 \\
\hline %za vodoravno crto
\multicolumn{2}{|c|}{\textbf{Povprečje}} & \multicolumn{2}{c|}{16} & 32 \\
\hline %\multicolumn za zdruzevanje stolpcevS
\end{tabular}
\end{table}


\begin{figure}[h] %[h] za here, %[t] za top, [b] za bottom, [p] za posebno stran za slike
\centering
\caption{muca}
\label{muca}
\includegraphics[width=0.5\textwidth]{muca.jpg}
\end{figure}

Zdaj pride pikazni način.
\[d\sin(ax)+b=c\] %prikazni nacin
\break
Ce se hoces sklicati na enacbo:
\begin{equation}\label{eq:pitagora}
    a^2 + b^2 = c^2
\end{equation}
Ali pa ce se hoces sklicati na enacbo ki nima stevilke zraven:
\begin{equation*}\label{eq:euler}
    e^{i\pi}+1=0
\end{equation*}
Potem pa še vrstični način.
$\frac{3x}{2y}-\sqrt{xy}=e^{2i}$ %vrsticni nacin

To: $\lim_{a \to x}{\ln{a}}$ je limita.

Se vsota:
\[\sum_{i=0}^{\infty}{\binom{i+2}{i}}\]

Se integral:
\[\int_{-\pi}^{\pi}{\cos^2(x)}dx\]

$ \begin{bmatrix}
\ddots & a_2 & a_3 & a_4 & a_5 & a_6 \\
b_1 & \ddots & b_3 & \ldots & \ddots & b_6 \\
c_1 & c_2 & 1 & 1 & \cdots & \ddots \\
d_1 & d_2 & d_3 & \vdots & d_5 & d_6 \\
e_1 & e_2 & e_3 & e_4 & e_5 & e_6 \\
f_1 & f_2 & f_3 & f_4 & f_5 & f_6 
\end{bmatrix}  $

\begin{multline*} %za enacbe daljse kot je sirina strani
p(x) = 3x^6 + 14x^5y + 590x^4y^2 + 19x^3y^3\\ 
- 12x^2y^4 - 12xy^5 + 2y^6 - a^3b^3
\end{multline*}

\begin{align*} %za poravnavo enacbo (&)
2x - 5y &=  8 \\ 
3x + 9y &=  -12
\end{align*}

Zdaj pa sklicevanje... \bigbreak %prehod v novo vrstico
Poglej sliko \ref{muca}
\bigbreak
Poglej tabelo \ref{tab:rezultati}
\bigbreak
Poglej enačbo \eqref{eq:pitagora}
\bigbreak
Poglej enačbo \eqref{eq:euler}
\bigbreak
Poglej stran \pageref{eq:pitagora}
\bigbreak
%\cite{vir} za citiranje ko imas bibliografijo
%\nocite{vir} ce citiras nekoga v bibliografiji ampak ga ne omenis v besedilu
%\bibliograpfystyle{plain}
%\bibliography{vir} za datoteko z bibliografijo
\end{document}
